\documentclass[11pt,a4paper]{article}

% ============================================================
% PAQUETES
% ============================================================
\usepackage[utf8]{inputenc}
\usepackage[T1]{fontenc}
\usepackage[spanish,es-tabla]{babel}
\usepackage[margin=2.5cm]{geometry}
\usepackage[table]{xcolor}
\usepackage{graphicx}
\usepackage{float}
\usepackage{tabularx}
\usepackage{array}
\usepackage{hyperref}
\usepackage{caption}
\usepackage{sectsty}
\usepackage{enumitem}

% ============================================================
% COLORES
% ============================================================
\definecolor{darkblue}{RGB}{0, 51, 102}
\definecolor{lightblue}{RGB}{204, 229, 255}

% ============================================================
% CONFIGURACIÓN DE HIPERVÍNCULOS
% ============================================================
\hypersetup{
    colorlinks=true,
    linkcolor=darkblue,
    filecolor=magenta,
    urlcolor=darkblue
}

% ============================================================
% ESTILO DE TÍTULOS
% ============================================================
\sectionfont{\color{darkblue}}
\subsectionfont{\color{darkblue}}
\subsubsectionfont{\color{darkblue}}

% ============================================================
% COMANDOS PERSONALIZADOS
% ============================================================
\newcommand{\code}[1]{\texttt{#1}}

% Comando para insertar imágenes con placeholder si no existen
\newcommand{\insertimage}[4]{
    \begin{figure}[H]
        \centering
        \IfFileExists{#1}{
            \includegraphics[
                width=0.95\textwidth,
                height=0.78\textheight,
                keepaspectratio
            ]{#1}
        }{
            \fbox{
                \begin{minipage}[c][7cm][c]{0.9\textwidth}
                    \centering
                    \Large \textbf{Imagen no encontrada}\\[0.5cm]
                    \normalsize #2\\[0.5cm]
                    \small Ruta esperada: \texttt{#1}
                \end{minipage}
            }
        }
        \caption{#3}
        \label{#4}
    \end{figure}
}

% Comando para insertar cuadros capturados como imagen
\newcommand{\insertcuadro}[3]{
    \begin{figure}[H]
        \centering
        \IfFileExists{#1}{
            \includegraphics[
                width=0.98\textwidth,
                height=0.82\textheight,
                keepaspectratio
            ]{#1}
        }{
            \fbox{
                \begin{minipage}[c][7cm][c]{0.9\textwidth}
                    \centering
                    \Large \textbf{Cuadro no encontrado}\\[0.5cm]
                    \small Ruta esperada: \texttt{#1}
                \end{minipage}
            }
        }
        \caption{#2}
        \label{#3}
    \end{figure}
}

% Comando para insertar capturas de código
\newcommand{\insertcodeimage}[3]{
    \begin{figure}[H]
        \centering
        \IfFileExists{#1}{
            \includegraphics[
                width=0.98\textwidth,
                height=0.82\textheight,
                keepaspectratio
            ]{#1}
        }{
            \fbox{
                \begin{minipage}[c][7cm][c]{0.9\textwidth}
                    \centering
                    \Large \textbf{Código no encontrado}\\[0.5cm]
                    \small Ruta esperada: \texttt{#1}
                \end{minipage}
            }
        }
        \caption{#2}
        \label{#3}
    \end{figure}
}

% ============================================================
% DOCUMENTO
% ============================================================
\begin{document}

\begin{center}
    {\Huge \color{darkblue} \textbf{Trabajo del Integrante 6: Frontend, API REST, pruebas e integración}} \\[1cm]
\end{center}

% ============================================================
\section{Datos Generales}

\insertcuadro
{imagenes/cuadro1.png}
{Datos generales del trabajo del integrante 6}
{cuadro:datos_generales}

% ============================================================
\section{Objetivo de la Parte 6}

El objetivo principal es implementar y evidenciar el desacoplamiento total entre el frontend y el backend mediante el uso de una API REST. Esto implica configurar la integración para el consumo de datos reales, sustituyendo los datos quemados o simulados, validar los endpoints generados, comprobar su comportamiento mediante pruebas y capturar de forma verificable el funcionamiento del sistema en conjunto.

% ============================================================
\section{Relación con las partes anteriores}

Este trabajo funge como la capa de consolidación e integración de todo el proyecto:

\begin{itemize}[leftmargin=1.2cm]
    \item Se apoya de forma directa en el \textbf{modelo de dominio} diseñado por el integrante 2, respetando las entidades de \code{Project} y \code{Task}.
    
    \item Consume activamente los módulos refactorizados de \code{Project}, correspondiente al integrante 3, y \code{Task}, correspondiente al integrante 4, interactuando con los controladores expuestos.
    
    \item Se alinea con los \textbf{Bounded Contexts} propuestos por el integrante 5, ya que la comunicación desde el frontend se realiza mediante peticiones discretas a la API REST, manteniendo los límites de contexto. En este caso, \code{projectId} actúa como puente relacional sin acoplar las lógicas internas.
\end{itemize}

% ============================================================
\section{Arquitectura de Integración}

La arquitectura se basa en una separación limpia cliente-servidor comunicada a través de un protocolo sin estado, HTTP. El frontend gestiona únicamente la capa de presentación y su estado interno, delegando cualquier lógica de negocio y persistencia de dominio al backend.

\insertimage
{imagenes/0.png}
{Figura principal de arquitectura}
{Arquitectura de Integración Frontend-Backend mediante REST}
{fig:arquitectura}

% ============================================================
\section{Desacoplamiento del Frontend}

Para evitar un diseño monolítico y un alto acoplamiento a las rutas de red, se implementó una \textbf{capa de abstracción de red}. Esta capa permite centralizar las peticiones realizadas desde React hacia el backend.

Se crearon y adaptaron los siguientes archivos para centralizar la comunicación:

\begin{itemize}[leftmargin=1.2cm]
    \item \code{api.js}: instancia principal de Axios configurada por la variable de entorno \code{REACT\_APP\_API\_BASE\_URL}. Centraliza cabeceras y permite fácil inyección de JWT en el futuro.
    
    \item \code{projects.js}: contiene los métodos para las llamadas a la API referentes a proyectos.
    
    \item \code{tasks.js}: contiene los métodos correspondientes a tareas.
\end{itemize}

\subsection{Ventajas del desacoplamiento}

\begin{itemize}[leftmargin=1.2cm]
    \item \textbf{Mantenibilidad:} si la IP o puerto del servidor cambia, únicamente se actualiza el archivo \code{.env}, no cada componente de la interfaz.
    
    \item \textbf{Reutilización:} distintos componentes pueden consumir el mismo servicio sin duplicar lógica de conexión.
    
    \item \textbf{Tolerancia al cambio:} facilita migrar hacia una arquitectura basada en microservicios en el futuro, cambiando algunas URLs en lugar de modificar toda la capa visual.
\end{itemize}

% ============================================================
\section{Endpoints REST Validados}

\insertcuadro
{imagenes/cuadro2.png}
{Endpoints REST validados en el sistema}
{cuadro:endpoints_rest}

% ============================================================
\section{Cambios principales en Frontend}

\insertcuadro
{imagenes/cuadro3.png}
{Cambios principales realizados en el frontend}
{cuadro:cambios_frontend}

% ============================================================
\section{Corrección de Deuda Técnica detectada por SonarQube}

El uso estático e inseguro de colecciones y componentes sin validar parámetros eleva los índices de \textit{Code Smells} y \textit{Reliability Issues} reportados por herramientas como SonarQube. Por ello, se aplicaron mejoras enfocadas en validaciones, control de valores nulos y correcciones visuales.

\insertcuadro
{imagenes/cuadro4.png}
{Corrección de deuda técnica detectada por SonarQube}
{cuadro:deuda_tecnica}

% ============================================================
\section{Pruebas Realizadas}

\subsection{Pruebas Backend}

Se ejecutó una compilación completa para validar la cobertura de pruebas REST en el entorno Spring Boot.

\begin{itemize}[leftmargin=1.2cm]
    \item \code{mvn clean test}: permitió validar el conjunto de pruebas usando MockMvc en \code{ProjectControllerTest} y \code{TaskControllerTest}.
    
    \item \code{mvn clean install}: finalizó con \textbf{BUILD SUCCESS}, validando el esquema de datos y generando los reportes correspondientes en el proyecto.
\end{itemize}

\subsection{Pruebas Frontend}

La interfaz fue validada unitariamente y reempaquetada para garantizar que no se introdujeran regresiones.

\begin{itemize}[leftmargin=1.2cm]
    \item \code{npm test -- --watchAll=false}: permitió ejecutar la suite de pruebas de React.
    
    \item \code{npm run build}: finalizó la compilación para producción, garantizando la integridad del empaquetado.
\end{itemize}

\subsection{Pruebas Manuales}

Se comprobó el paso de mensajes de extremo a extremo:

\begin{itemize}[leftmargin=1.2cm]
    \item Navegación completa mediante React, asegurando que las acciones de crear, editar y borrar se reflejen tras recargar la página.
    
    \item Ejecución manual directa desde un archivo \code{api-tests.http} en VS Code sobre los endpoints de proyectos y tareas.
\end{itemize}

% ============================================================
\section{Evidencias Funcionales}

\insertcuadro
{imagenes/cuadro5.png}
{Relación de evidencias funcionales del sistema}
{cuadro:evidencias_funcionales}

\insertimage
{imagenes/backend-ejecucion.png}
{Terminal con Spring Boot}
{Backend ejecutándose en puerto 8080}
{fig:backend_run}

\insertimage
{imagenes/frontend-ejecucion.png}
{Navegador en localhost:3000}
{Frontend levantado e interfaz activa}
{fig:frontend_run}

\insertimage
{imagenes/postman-get-projects.png}
{Petición GET de proyectos}
{Prueba GET proyectos}
{fig:postman_get_projects}

\insertimage
{imagenes/postman-post-project.png}
{Petición POST de proyecto}
{Prueba POST proyecto}
{fig:postman_post_project}

\insertimage
{imagenes/postman-get-tasks.png}
{Petición GET de tareas}
{Prueba GET tareas por proyecto}
{fig:postman_get_tasks}

\insertimage
{imagenes/react-lista.png}
{Listado de Projects y Tasks en React}
{Vista React consumiendo la API REST}
{fig:react_consuming_api}

\insertimage
{imagenes/react-crea-proyecto.png}
{Formulario React para proyecto}
{Creación de proyecto desde frontend}
{fig:react_create_project}

\insertimage
{imagenes/react-crea-tarea.png}
{Formulario React para tarea}
{Creación de tarea desde frontend}
{fig:react_create_task}

\insertimage
{imagenes/react-build.png}
{Terminal con comando npm run build}
{Build optimizado y exitoso del frontend}
{fig:react_build}

\insertimage
{imagenes/backend-build-success.png}
{Terminal con BUILD SUCCESS de Maven}
{Pruebas y build exitoso del backend}
{fig:backend_tests}

% ============================================================
\section{Fragmentos de Código Relevantes}

En esta sección se presentan los fragmentos de código más importantes empleados en la integración entre el frontend y el backend.

\subsection{Configuración dinámica de API}

\insertcodeimage
{imagenes/code1.png}
{Configuración dinámica de API en api.js}
{code:api_js}

\subsection{Implementación de PropTypes de defensa en listas}

\insertcodeimage
{imagenes/code2.png}
{Implementación de PropTypes en TasksList.js}
{code:taskslist_js}

\subsection{Validación de Backend consumida por la API}

\insertcodeimage
{imagenes/code3.png}
{Validación de backend en ProjectController.java}
{code:projectcontroller_java}

\subsection{Archivo de prueba REST automatizable}

\insertcodeimage
{imagenes/code4.png}
{Archivo de prueba REST api-tests.http}
{code:api_tests_http}

% ============================================================
\section{Resultados Obtenidos}

Gracias a los pasos de integración y refactorización aplicados, se alcanzaron los siguientes logros clave:

\begin{itemize}[leftmargin=1.2cm]
    \item \textbf{Frontend totalmente desacoplado:} la interfaz de usuario opera bajo un esquema cliente-servidor puro. No existen datos quemados y todas las operaciones se envían mediante peticiones HTTP.
    
    \item \textbf{Comunicación fluida mediante API REST:} la integración funcionó correctamente tras configurar la comunicación entre React y Spring Boot.
    
    \item \textbf{Aumento de confiabilidad:} las adiciones de \code{PropTypes} y protecciones contra estados inconsistentes lograron corregir defectos identificados y proteger el código contra deuda técnica.
    
    \item \textbf{Verificación general integral:} la ejecución satisfactoria de comandos como \code{mvn install} y \code{npm test} demuestra un estado sano del repositorio.
\end{itemize}

% ============================================================
\section{Conclusiones}

\begin{enumerate}[leftmargin=1.2cm]
    \item La separación en capas impuesta por DDD facilita la integración con un frontend independiente porque cada contrato entre sistemas queda definido desde el controlador mediante DTOs.
    
    \item El uso de variables de entorno, como el archivo \code{.env} en React, es indispensable para el ciclo de desarrollo ágil, ya que evita codificar rutas o puertos directamente en múltiples archivos.
    
    \item La implementación de bibliotecas como \code{prop-types} en React disminuye errores de ejecución y mejora la confiabilidad de los componentes.
    
    \item Al definir las relaciones únicamente con identificadores de contexto, como \code{projectId}, se logra una cohesión alta y un acoplamiento bajo.
    
    \item Las métricas de calidad de código no solo afectan el backend; asegurar contrastes de color apropiados y protección de arreglos en el frontend también es clave para el mantenimiento del sistema.
\end{enumerate}

% ============================================================
\section{Resumen del Aporte}

\insertcuadro
{imagenes/cuadro6.png}
{Resumen del aporte realizado por el integrante 6}
{cuadro:resumen_aporte}

\end{document}